\documentclass[12pt,a4paper]{ctexart}
\usepackage{graphicx}
\usepackage{booktabs}
\usepackage{float}
\usepackage{geometry}
\geometry{margin=2.5cm}
\usepackage{caption}
\usepackage{subcaption}
\usepackage{hyperref}

\begin{document}

% 封面
\begin{titlepage}
    \centering
    {\Huge\bfseries 空中机器人技术方案\par}
    \vspace{1cm}
    {\large\bfseries RoboMaster 2025\par}
    \vspace{1cm}
    {\large 组织：上海交通大学 云汉交龙\par}
    \vspace{0.5cm}
    {\large 时间：2025年8月\par}
\end{titlepage}

\tableofcontents
\newpage

\section{前言}
本技术报告由上海交通大学云汉交龙战队编制，适用于RoboMaster 2025机甲大师超级对抗赛。主要撰写人员包括：

\begin{table}[h]
    \centering
    \begin{tabular}{lll}
        \toprule
        模块 & 主要撰写人员 & 协助撰写人员 \\
        \midrule
        机械 & 郭锦鸿 & N/A \\
        \bottomrule
    \end{tabular}
    \caption{撰写人员表}
\end{table}

\section{背景与目标}
在RoboMaster比赛中，空中机器人的作用日益凸显。在25赛季新规则改动中，无人机允许发弹量增加到1500发且无热量上限，这极大的增加了空中机器人在赛场上的压制力。无人机在赛场上能干的事情非常多：能快速击打前哨站、击杀敌方地面关键单位、攻打敌方基地等。而且得益于25赛季地形改动，使得空中机器人激活能量机关成为可能，所以无人机还能开符。总而言之优秀的空中机器人是一个全能型选手，既能进攻建筑物还能远距离击杀敌方关键单位，既能冲锋进攻也能守家，甚至还能开符获得团队增益效果。

列出以下几条具体要求与目标：
\begin{itemize}
    \item 稳定的飞行能力：包含性能强劲的动力系统与抗干扰能力优秀的飞控系统
    \item 定点定位系统：感知自身飞行位置与姿态，保持定点稳定飞行
    \item 充沛的续航时间：在实战场景下轻松飞行7min以上
    \item 轻量化高强度机身结构：机身整体刚性强且质量轻
    \item 精确的发射系统：高频顺畅发射不卡弹，散度10m一块小装甲
    \item 大角度双轴云台：给云台手获得宽广视野、能近距离开符
\end{itemize}

\section{机械结构设计}
无人机的整体机械结构在大体上可以分为飞行平台与发射系统两个部分，每个部分细分可以具体分为机身、机臂、脚架、云台等结构，详细关系见下图。

% 图片占位：图2.1
\begin{figure}[H]
    \centering
    % 请替换为实际图片文件，如 fig2_1.png
    \includegraphics[width=0.8\textwidth]{fig2_1.png}
    \caption{飞行平台结构}
    \label{fig:structure}
\end{figure}

\subsection{飞行平台机械结构设计}
飞行是一台无人机最基础的功能，所以飞行平台的结构是无人机最重要的结构。对它的要求是：
\begin{enumerate}
    \item 整体刚度大，能够减少无人机在飞行过程中的震动，并且在发生意外坠机时能保证不发生严重破坏。
    \item 质量轻，轻量化设计理念会给终身穿整个无人机机械设计过程，而飞行平台是减重的大头。
\end{enumerate}

无人机在飞行时会产生很多高频震动，如果机身设计不稳定或刚度不够会极大程度影响飞行平稳性和飞机的操控性，进而影响发射系统，所以优秀的飞行平台设计是赛场上空中机器人发挥优秀的基础。接下来依次介绍飞行平台各个部分的设计以及之间的联系。

% 图片占位：图2.2
\begin{figure}[H]
    \centering
    % 请替换为实际图片文件，如 fig2_2.png
    \includegraphics[width=0.8\textwidth]{fig2_2.png}
    \caption{飞行平台装配}
\end{figure}

\subsubsection{机身中心结构}
机身的中心结构是整个飞行平台的中心，所有其他的结构都固定在机身中心板上，向四周伸出四根机臂，下面接脚架，上面接弹舱，而且飞控、主控等电子元件也固定在其中。它的组成结构很简单，就是由上下两层碳板中间加铝柱构成，如图2.2所示。

% 图片占位：图2.3
\begin{figure}[H]
    \centering
    % 请替换为实际图片文件，如 fig2_3.png
    \includegraphics[width=0.8\textwidth]{fig2_3.png}
    \caption{机身中心结构}
\end{figure}

因为要与机身其他机构连接所以做了很多固定的孔位，并在空白处做大量镂空减重。所以从中心板的孔位布局能够看出飞行平台的布局方式。以中心下板为例说明孔位功能：

% 图片占位：图2.4
\begin{figure}[H]
    \centering
    % 请替换为实际图片文件，如 fig2_4.png
    \includegraphics[width=0.8\textwidth]{fig2_4.png}
    \caption{下中心版孔位}
\end{figure}

蓝色圈1是机臂方管的固定孔，红圈2是脚架的固定孔，黄圈3是给分电板吊装预留的，绿圈4是云台架的固定孔位，5是留给弹舱波纹管的通道，周边是波纹管固定孔，其他孔位是上下两层中心板之间的加强铝柱孔。所以机身中心结构的设计重点是考虑整个飞行平台的布局，先把整机尺寸与各个部分的相对位置确定下来以后，中心机构自然就设计出来了。建议首先从飞机的功能出发，确定机臂大致尺寸、云台是否前置、脚架位置等因素，才能够把整机的布局确定。总的来说机身中心是整个飞行平台的固定点，所以中心板需要高强度的同时各元件布局也需要合理且紧凑。

\subsubsection{机臂设计}
25赛季无人机机臂相比去年有很大提升，由之前的圆管换成薄壁碳方管带来了许多优点。

% 图片占位：图2.5
\begin{figure}[H]
    \centering
    % 请替换为实际图片文件，如 fig2_5.png
    \includegraphics[width=0.8\textwidth]{fig2_5.png}
    \caption{机臂装配体}
\end{figure}

因为旋翼电机、桨叶等需要直接与机臂相连，在飞行时需要承受很大拉力需要一定的强度，但同时也需要保持桨叶平面水平。而方管形变比圆管小，而且方管不需要调节桨平此外还能在侧面镂空减重，所以今年选用的是 $30^{\circ}30$ 壁厚 $1\mathrm{mm}$ 的碳方管。在机臂上安装电机电调等比较简单，直接钉上去即可，不太需要额外增加结构与重量。单根方管机臂相比去年的圆管机臂轻 $50\mathrm{g}$，而整个机臂装配体相比去年轻将近 $800\mathrm{g}$。

方管机臂侧面的镂空采用拓扑方法，通过Altair inspire软件拓扑仿真后得出具体镂空形状，如图2.5所示。

% 图片占位：图2.6
\begin{figure}[H]
    \centering
    % 请替换为实际图片文件，如 fig2_6.png
    \includegraphics[width=0.8\textwidth]{fig2_6.png}
    \caption{方管机臂拓扑优化}
\end{figure}

\subsubsection{脚架结构设计}
无人机脚架的作用是支撑整架飞机，在迫降甚至坠机时是受力最大的结构，而且脚架离机身最远，最容易产生震动从而影响飞行效果，所以需要极高的稳定性与刚度。

脚架尺寸的最大的影响因素是云台，因为云台在脚架中间运动时很容易和脚架发生干涉，所以先定云台的活动范围，再定脚架的长度以及间距。由于25赛季新增打符功能，云台仰角和俯角都很大所以把脚架加长到 $550\mathrm{mm}$。并且为了获得更大视野，云台yaw轴活动范围可达到 $180^{\circ}$ 还做了向前偏置，脚架间距也因此调整，需要云台保证任意角度不与脚架干涉。

脚架与机身中心下板连接处采用了快拆管夹转接件，优点是方便拆卸，能快速处理脚架损坏，但由于脚架尺寸过于长所以管夹处容易出现松动或晃动，需要加碳纤维止滑剂并用维拉扳手拧死。脚架下端加了横杆约束脚架的横向位移，形成更加稳定的三角形结构，因为脚架尺寸进行了加长了所以很需要这个结构。赛场上表现不错，在轻度坠机情况下仍然能保证脚架主体结构的完整性，横杆功不可没。

% 图片占位：图2.7
\begin{figure}[H]
    \centering
    % 请替换为实际图片文件，如 fig2_7.png
    \includegraphics[width=0.8\textwidth]{fig2_7.png}
    \caption{脚架结构}
\end{figure}

\subsubsection{弹舱设计}
弹舱设计首先考虑容量，比赛中允许最大发单量1500发，弹舱容积最好在9L以上，但最需要考虑的是是否会和其他结构干涉。因为采用上供弹结构的发射系统，所以弹舱一般都在上中心板上方，可能会与电池、gps平台等元件干涉，而且弹舱的尺寸和形状在很大程度上取决于上中心板的尺寸和电池的摆放位置，所以弹舱一般最后设计。很多学校的弹舱采用打印件甚至塑料等，但我们用的1mm碳板强度会更大，赛场表现还不错。

\subsection{动力系统}
25赛季的动力系统采用的是好盈 8100UL 旋翼电机 + 好盈 3098 桨叶，这套动力系统在25赛季中使用非常广泛，很多学校都采用这套配置，而且在本赛季规则下有亮眼表现。满载整机重量将近 $18\mathrm{Kg}$ 但这套动力系统毫无压力，而且温度正常，从没出现过热等异常情况。

30寸的大桨叶提升气动效率，增加续航时间而且更大的桨叶能产生更大的总体升力。同时还能改善飞行稳定性和悬停性能，大桨叶具有更大的转动惯量，就像陀螺仪一样，有助于稳定无人机的姿态。大桨叶在低转速下工作，产生的气流更平缓，湍流更少，这使得无人机在悬停和低速飞行时更加平稳。但是大桨叶带来的缺点是让整架无人机更难控制，任何微小的不平衡都会被放大，导致整机产生更剧烈的振动，严重影响飞控传感器的精度，所以对机身结构的稳定性要求更高。

\subsection{桨保持结构设计}
25赛季采用的是一体化桨保，整个飞机看起来像一个方盒子。最初考虑做一体式桨保的动机是想减重，四个分体式桨保的重量加起来有 $1.6\mathrm{Kg}$，但一体式桨保大概 $1.2\mathrm{kg}$，而且还有一个非常重要的点是一体式桨保的刚度很高，对内部结构能起到有效的保护，在搬运过程中也能感受出整个一体式桨保很坚固。但是最大的缺点是不方便维护，飞机内部的中心板和飞控都被桨保覆盖，每次拔插飞控线或者检修内部结构的时候都需要拆掉桨保，虽然已经尽量做了三层式快拆设计但仍然很不方便，建议可以把飞控等需要经常维护的元件放在面上，别让桨保覆盖住了。桨保也很容易与电池发射干涉，因为电池需要方便拔插，肯定不能被桨保覆盖，所以25国赛飞机一改往日把电池放机身里的传统，索性设计了新的电池架结构直接放脚架上，效果也挺好。

覆盖网选用的是 6 股 $2.5\mathrm{cm}$ 网眼的渔网，这个渔网的质量最轻，虽然并不是很结实但也够用，能满足规则手册中的检录要求。它并不能够防止小弹丸进入桨叶但30寸的大桨叶不怕小弹丸，不会影响飞行姿态，但是可能会把桨叶侧边打裂发生分层。

\subsection{发射系统}
RM可以看作一个发射比赛，空中机器人的所有的攻击都是依靠发射小弹丸造成伤害，所以飞行平台是空中机器人的基础，发射系统能提升它的上限，优秀的无人机一定有优秀的发射。

我们采用的是上供弹结构，弹丸从弹舱中流向拨盘，然后在拨盘的作用下一颗颗进入链路中，最后通过一对摩擦轮发射出去。在机械上要做到云台活动范围大，任意角度高频发射不卡弹，发射精度高散度低。

% 图片占位：图2.8
\begin{figure}[H]
    \centering
    % 请替换为实际图片文件，如 fig2_8.png
    \includegraphics[width=0.8\textwidth]{fig2_8.png}
    \caption{一体式快拆架保}
\end{figure}

% 图片占位：图2.9
\begin{figure}[H]
    \centering
    % 请替换为实际图片文件，如 fig2_9.png
    \includegraphics[width=0.8\textwidth]{fig2_9.png}
    \caption{发射系统}
\end{figure}

\subsubsection{云台架与拨盘}
云台架是连接机身与云台的部分，整个云台通过它吊装固定在机身下中心板上。由于25赛季我们的无人机新增了开符功能，云台增加了向上 $25^{\circ}$ 的仰角，为了防止弹丸发射弹道与一体式桨保干涉所以云台架变得特别长，我们叫它“长脖子”。

说到“长脖子”就不得不提我们曲折的设计过程。为了获得仰角不得不把带有发射机构的云台下降，最开始我们只是单纯的加长竖直段的链路，拨盘还是像24赛季无人机一样留在中心板上，然后云台用四根粗铝柱吊装，如下图所示。这个方案有许多问题，首先就是链路长度太长，中间蓝色的打印件就是一段非常长竖直段链路，长度来到了 $230\mathrm{mm}$，带来的问题是弹丸阻力太大，拨盘根本拨不动。我们尝试过很多减小竖直段链路阻力的办法，例如加小轴承或将圆柱拆分成四片但都效果不佳，链路阻力依旧很大。这个方案第二个致命问题是整个吊装云台可以看成是很长的悬臂梁，云台运动产生的扰动会通过悬臂梁放大，导致整个云台特别晃，根本无法稳定运行，而且起飞后对飞行的整体扰动也会很大。

% 图片占位：图2.10
\begin{figure}[H]
    \centering
    % 请替换为实际图片文件，如 fig2_10.png
    \includegraphics[width=0.8\textwidth]{fig2_10.png}
    \caption{初始版云台吊装方案与云台架吊装方案}
\end{figure}

要解决上述两个问题，第一必须减小链路长度，第二需要加固云台吊装方式，所以最终版云台架被我们慢慢摸索出来，如上图所示。

为了缩短链路长度将拨盘下沉，把原本在中心板内的拨盘下降到云台上方，用单独一层碳板固定拨盘与直角链路，再用波纹管连接拨盘与弹舱，这样消除了中间段 $230\mathrm{mm}$ 的竖直段链路。然后用三块厚碳板形成方形结构来吊装云台，能抵抗扭转与横向位移，大大增加了云台的稳定性，有效减小因为云台运动产生的晃动与形变。

拨盘用的是2006电机老拨盘，虽然有弹丸卡挡边等问题但性能也还算稳定，能满足发射需求。

\subsubsection{云台链路}
云台的链路设计部分首先需要明确云台pitch轴的仰角与俯角，然后根据链路设计守则画草图，链路的转弯半径 $24\mathrm{mm}$，再设计对应的活动挡边，这部分和队里传统一样，关键是画好草图，否则就会在某些角度阻力增大导致弹丸运行不顺畅。

25国赛空中机器人采用双轴云台，yaw轴可活动范围 $180^{\circ}$，pitch轴仰角 $25^{\circ}$ 俯角 $45^{\circ}$，这个活动范围可以算是非常大了，在比赛中能给云台手提供非常大的视野，能够在固定点拥有广阔的打击范围而不用频繁调整飞机位置。仰角 $25^{\circ}$ 能够实现近距离击打能量机关，这在开大符的时候非常有用，近距离击打能够有效提升命中环数。我们这次的空中机器人在国赛中表现亮眼，打出本赛季最高39环的优秀成绩。

% 图片占位：图2.11
\begin{figure}[H]
    \centering
    % 请替换为实际图片文件，如 fig2_11.png
    \includegraphics[width=0.8\textwidth]{fig2_11.png}
    \caption{云台链路}
\end{figure}

有一点需要注意的是pitch轴采用四边形连杆传动，压力角不能太小否则电机带不动，而且尽量设计成平行四边形连杆传动，传力更均匀。

\subsubsection{云台原件布局}
在云台下部是一块大碳板，各个元件都布置在它上面。这些元件有电调、电管、小电脑等，主要是与发射相关的电子部件。这部分是线路最复杂的地方，所以25赛季我们在云台上增加了一块PCB板来替代杂乱的线束，使得整个云台看起来清爽很多，而且方便修理维护。因为云台上的元件数量多，很容易在装配是发生干涉所以每个元件的布局都需要考虑清楚。特别是PCB板与电管、电调的固定，可以想一想装配的顺序，不然很容易出现无法装配的屎。

% 图片占位：图2.12
\begin{figure}[H]
    \centering
    % 请替换为实际图片文件，如 fig2_12.png
    \includegraphics[width=0.8\textwidth]{fig2_12.png}
    \caption{云台元件与发射机构}
\end{figure}

我们将云台的位置前置，并不在传统的飞机中心点上，目的是为了获得更大yaw轴视野，而且为了在左右旋转时不与脚架发生干涉特意缩短了云台板的长度，最终实现云台水平 $180^{\circ}$ 旋转，非常有实战意义。

为了自瞄能锁定远处的目标与近处的能量机关，设计了双相机模块，采用 $12\mathrm{mm}$ 与6mm焦距的相机来应付远近不同的目标。在激活能量机关时会启用6mm焦距镜头，其他时间都使用 $12\mathrm{mm}$ 焦距镜头锁定目标。双相机安装件在左右两边还设有c板和imu的安装孔位，能更方便地接线与维护。而且该模块做了快拆设计，把整个加工件分为三部分，这样能在视觉标定时方便拆卸，后来这部分确实拆了很多次，得益于快拆设计并没有很费工夫。

% 图片占位：图2.13
\begin{figure}[H]
    \centering
    % 请替换为实际图片文件，如 fig2_13.png
    \includegraphics[width=0.8\textwidth]{fig2_13.png}
    \caption{双相机模块}
\end{figure}

\subsubsection{发射机构}
17mm小弹丸是依靠一对摩擦轮的旋转发射出去，而单发限位是发射机构中非常重要的结构，往年都是采用地面车组的拉簧式单发限位，但拉簧对弹道的影响很大，非常影响发射散度。今年采用轴承式单发限位，通过在起始件上加一对U型轴承来限制预制位弹丸的活动，因为轴承是固定的，不会像拉簧一样反复入侵链路中，所以不会影响发射散度，而且轴承能够使每个出去的弹丸都在固定位置起到定心效果甚至会提升发射精度，这也是今年无人机发射散度小的秘诀。

% 图片占位：图2.14
\begin{figure}[H]
    \centering
    % 请替换为实际图片文件，如 fig2_14.png
    \includegraphics[width=0.8\textwidth]{fig2_14.png}
    \caption{轴承单发限位起始件}
\end{figure}

轴承型号如下图所示，轴承间距是关键参数，通过做demo反复测试得出 $28.6\mathrm{mm}$ 是最合适的间距，如果间距太大弹丸卡不住，起不到单发限位效果，如果太小则阻力太大，拨盘不能顺畅转动。轴承的安装位置在摩擦轮之后，注意不要和摩擦轮发射干涉。

% 图片占位：图2.15
\begin{figure}[H]
    \centering
    % 请替换为实际图片文件，如 fig2_15.png
    \includegraphics[width=0.8\textwidth]{fig2_15.png}
    \caption{轴承尺寸参数}
\end{figure}

\appendix
\section{国赛空中机器人检赛前一小时与场间三分钟}

\subsection*{赛前一小时：}
\begin{enumerate}
    \item 自瞄：手持装甲板检测，连接小电脑看自瞄情况
    \item 定位：t265连接是否稳定
    \item 打弹：检查弹速与卡弹情况
    \item 检查线：是否都扎好、接头是否松动
    \item 电池电压（用手机拍照分高低组）
    \item 遥控器通道映射
    \item 检查螺丝是否紧固（桨叶、云台、脖子等）螺丝检查：1，pitch轴电机螺丝2，云台连杆螺丝3，链路螺丝4，起始件螺丝5，相机总装螺丝6，桨叶螺丝7，弹舱螺丝8，gps平台与保护杆螺丝9，拨盘与脖子螺丝10，桨保螺丝
\end{enumerate}

\subsection*{场间三分钟：}
\begin{enumerate}
    \item 拔电池
    \item 装1500发小弹丸
    \item 检查电池是否装紧并上电
    \item 发射打火、确保发射机构正常
    \item 等待自瞄灯亮起
    \item 插拔 T265 并检查 gps 灯看是否正常工作
\end{enumerate}

\end{document}